
\needspace{4\baselineskip}
\color{uniblue}{\subsection{\raggedright{Управление выключателем (УВ)}}}
\color{black}

\begin{enumerate}

%%%%%%%%%%%%%%%%%%%%%%%%%%%%%%%%%%%%%%%%%%%%%%%%%%%%%%%%%%%%%%%%%%%%%%  ОБЩИЕ СВЕДЕНИЯ %%%%%%%%%%%%%%%%%%%%%%%%%%%%%%%%%%%%%%%%%%%%%%%%%%
\item
УВ предназначено для передачи команды на включение выключателя от контроллера присоединения, в случае, когда контроллер присоединения является внешним устройством. В устройстве управление выключателем представлено функцией «УВ». Функция контролирует условия разрешения для выполнения команды включения выключателя.

\item\label{sec:uv}
Уставкой SGF2 «Контроль\_синхр» регулируется включение выключателя с контролем синхронизма или без него. Если уставка SGF2 предусматривает включение с КС, то наличие дискретного сигнала оперативного включения выключателя из КП и сигнала разрешения оперативного включения из логики КС формируют команду включения выключателя. Команда включения выключателя сбрасывает триггер, восстанавливая схему. Сброс триггера также осуществляется при поступлении сигнала блокировки из логики КСВ, сформировавшегося сигнала задержки включения выключателя с КС по истечении выдержки времени таймера Т2 или выводом функции из работы.

\item
Если уставка SGF2 «Контроль\_синхр» не предусматривает выключение выключателя с КС, то команда включения выключателя сформируется, не учитывая сигнала из логики КС.

\item
Ввод функции в работу на этапе параметрирования устройства осуществляется программным переключателем SGF1 «Ввод\_функции». Введенное состояние функционального блока соответствует активному состоянию сигнала «УВ: Ввод».

\item
Функция формирует сигнал «УВ / УВ: Включить» при активации сигнала «Опер. включить В» (команда включения от внешнего КП), разрешения включения от «КС» (см. п. \ref{sec:uv}) и при отсутствии сигнала блокировки включения выключателя от функции контроля выключателя «КСВ». Сигнал оперативного отключения от внешнего КП не требует проверки каких-либо условий и передается напрямую в схему управления выключателем.

\item 
Функция может быть выведена из работы оперативно путем активации сигнала «ОВ УВ», о чем сигнализирует активное состояние сигнала «УВ: Оперативный вывод».

\item
Алгоритм работы функции выполнен в соответствии с рисунком \ref{uv:img_gen}. В таблице \ref{uv:tbl_gen} приведены параметры, необходимые для настройки функции. 

\vspace{1mm}
\begin{figure}[H]
\centering
\includegraphics[width=1\textwidth,height=1\textheight,keepaspectratio]{\fbpath/129. УВ ОЛ/_latex/img/uv.pdf}
\captionof{figure}{Функциональная схема логики УВ}\label{uv:img_gen}
\end{figure}

\end{enumerate}
